%% kapa: exam=1 round=19 year=2008 period=2 form=A subject=economics subject_ko=경제원론 seq=3 number=3
%% source: Q-Net 감정평가사 기출문제 (cst00309) artlSeq=1001962 file=extracted/1008/1001962/2008년 제19회 감정평가사 1차시험 기출문제/2교시.hwp
%% conversion: kordoc 4.24.51 -> tools/gen1.py (auto); review=proofread-sonnet; flags=none
\begin{question}{3}
어떤 상품의 공급곡선은 $Q_{ s } = -100+4P$ 이고 수요곡선은 $Q_{ d } =150-P$ 이다. $P$, $Q_{ s }$, $Q_{ d }$ 는 각각 가격, 공급량, 수요량을 나타낸다. 정부가 상품 1개 당 25원의 세금을 생산자에게 부과하는 경우 정부의 세금수입은 얼마인가?
\begin{choices}
  \item 1,600원
  \item 2,000원
  \item 2,500원
  \item 3,000원
  \item 3,500원
\end{choices}
\end{question}
\answer{2}
